\newpage
\section*{Заключение}
\addcontentsline{toc}{section}{Заключение}

В данной работе были рассмотрены точки Лагранжа в ограниченной круговой задаче
трёх тел. Было показано, что при переходе к вращающейся системе координат
возникают пять особых точек равновесия, в которых малое тело может сохранять
постоянное положение относительно двух основных тел. Эти точки делятся на
коллинеарные $L_1$, $L_2$, $L_3$ и треугольные $L_4$, $L_5$.

Было отмечено, что коллинеарные точки являются неустойчивыми, тогда как
треугольные точки при выполнении условия на массовый параметр могут быть
устойчивыми. Именно это различие определяет характер движения малого тела в
их окрестности.

В вычислительной части эти выводы были подтверждены численно на примерах
точек $L_1$ и $L_4$ для системы Солнце - Земля. При задании начальных условий
рядом с $L_1$ траектория удалялась от положения равновесия, тогда как при
старте вблизи $L_4$ движение оставалось ограниченным в окрестности точки.
Таким образом, исследование точек Лагранжа является хорошим примером того,
как методы теоретической механики и численного моделирования дополняют друг
друга.
